%% ─────────────────────────────────────────────────────────────────────────────
\section{Plan for Next Steps}
%% ─────────────────────────────────────────────────────────────────────────────

\noindent\textbf{Limitations to address.} (i) Only Human-written T0 has been analysed; the six
LLM sources are preprocessed and will be compared next, including whether a model changes its
own text less (OpenAI T0 under ChatGPT, PaLM T0 under PaLM2). (ii) Paraphrases are shorter than
their sources and short texts give noisier profiles, so we will add a length-matched baseline.
(iii) SBERT reads only 256 tokens, which truncates 43\% of documents; a chunked encoding is
implemented, and a first check raises SBERT at T3 by 0.01--0.05 in long-text domains (ChatGPT on
stories: 0.817$\to$0.861), so final results will use it.

\noindent\textbf{Weeks 9--11: Authorship attribution (RQ2).} Train a stylometric SVM and a
fine-tuned RoBERTa to identify the seven sources on the train files, evaluate them on T0--T3, and
locate the ``point of no return'' per paraphraser.

\noindent\textbf{Weeks 10--12: Paraphraser fingerprinting (RQ3).} Train a paraphraser classifier
on T1--T3 (excluding DIPPER-high) with a key-level split, and visualise identity trajectories with
t-SNE and quantitative cluster-separation measures.

\noindent\textbf{Weeks 13--14: Synthesis.} A Streamlit decay dashboard and the final paper.
